\documentclass[11pt]{article}
\usepackage[margin=1in]{geometry}
\usepackage{amsmath,amssymb}
\usepackage[hidelinks]{hyperref}
\title{Independence of Two Transport Conditions}
\author{}
\date{Companion note to Interactads}
\begin{document}
\maketitle
\noindent This note separates the two grade-transport tests used in the
comma presentation of Interactads: exact base change and preservation
of cartesian grade multiplication. Both witnesses use finite surjections,
identity unit and multiplication bodies and identity occurrence, and satisfy the display
test. They address necessity of checking the stated hypotheses; they do
not obstruct a different representation of mixed composition.

\section{Common grade model}\label{sec:grade}
For a functor $K:\mathbf{Set}\to\mathbf{Set}$, take grades over $B$ to be
subsets $Q\subseteq K(B)$. An arrow $(B,Q)\to(B',Q')$ over $u:B\to B'$
exists precisely when $K(u)[Q]\subseteq Q'$, and is unique if it exists.
The arrows representing direct and inverse image are opcartesian and
cartesian respectively: their factorization conditions are
$K(u)[Q]\subseteq Q'$ if and only if $Q\subseteq K(u)^{-1}[Q']$.
Choose the grade extension to be the underlying base itself, and take
the occurrence monad and occurrence map to be identities. The display
comparison is then the ordinary pullback of a base map and is invertible.

Equip each grade fibre with an associative unital multiplication that
direct image preserves. A content arrow consists of a body $f:A\to B$
and a grade $Q$ over $B$. Following $(Q,f)$ by $(R,g)$ uses grade
$R\cdot g_!Q$ and body $gf$. The multiplication, functoriality and unit
laws give the one-sided content calculus. All content unit and
multiplication bodies are identity maps. Horizontally accumulating the
chosen cartesian grade arrows along $u$ remains cartesian exactly when
$u^*(R\cdot Q)=u^*R\cdot u^*Q$; the cartesian unit test is preservation
of the unit grade under inverse image.

\section{Base change can fail with cartesian grade multiplication}\label{sec:bc}
Let $K(B)=\mathcal P_{\le3}(B)$, where $K(u)$ takes direct image.
Grade multiplication is union of collections of subsets, and the unit
is the empty collection. Both direct and inverse images preserve union
and the empty collection. Thus the cartesian algebra condition holds.

Take $X\to B\leftarrow Y$ where $B=\{0,1\}$, $X$ has fibre sizes $(2,1)$,
and $Y$ has fibre sizes $(1,2)$. Put $Z=X\times_B Y$ with projections
$p:Z\to X$ and $q:Z\to Y$. For the grade $Q=\{X\}\subseteq K(X)$,
the image of $X$ is $B$, so $Y$ belongs to the inverse image on $Y$ of
the direct image of $Q$. To belong to $q_!p^*Q$ it would have to be the
$q$-image of a subset of $Z$ of size at most three whose $p$-image is $X$.
Over $0$ two elements are needed to cover $X$; over $1$ two are
needed to cover $Y$. This needs four elements. Hence the canonical
Beck--Chevalley inclusion is strict. Both input arrows and both
projections are finite surjections.

\section{Grade multiplication can fail with exact base change}\label{sec:mult}
Let $K(B)=\mathsf{List}(B)\times\mathbb N$, with multiplication
\[
 (w,n)(v,m)=(w\,\operatorname{dup}^{n}(v),n+m),
\]
where $\operatorname{dup}$ duplicates each letter and juxtaposition
concatenates. This is a monoid by the action of $\mathbb N$ by list
endomorphisms. Direct images of subsets preserve the induced subset
product and its unit $\{([],0)\}$. The functor $K$ preserves pullbacks:
two lists with the same image have the same length, so their positions
pair uniquely, and the numerical coordinates match. In the preserved
pullback the two direct/inverse image routes therefore give the very
same subsets.

For the finite surjection $\pi:\{0,1\}\to\{*\}$ take grades
$R=\{([],1)\}$ and $Q=\{([*],0)\}$. The inverse image of $R\cdot Q$
contains $([0,1],1)$ because its image is $([*,*],1)$.
The product $\pi^*R\cdot\pi^*Q$ contains only $([0,0],1)$ and
$([1,1],1)$. The accumulated arrow thus fails the cartesian test: its
underlying grade is strictly smaller than the inverse image of the
accumulated grade. The unit still pulls back to the unit.

\end{document}
